% Unit 2 exam -- 20 MCQ (no calculator) + 1 long FRQ (10 pts). 52 min.
\documentclass{apchem-exam}

\apcunit{2}
\apcunittitle{Compound Structure and Properties}
\apcdoctitle{Unit Exam}
\apcsubtitle{Wed Sep 9 \textbullet\ 52 minutes \textbullet\ periodic table provided}

\begin{document}
\apctitleblock

\examsection{I}{Multiple Choice}{20 questions \textbullet\ 30 minutes \textbullet\ 1 point each}{\nocalculator}

% ---- 2.1 bond types --------------------------------------------------------
\begin{mcq}
  Which bond is the most polar?
  \choices
    \choice{\ce{C-H}}
    \choice{\ce{N-H}}
    \correct{\ce{O-H}}
    \choice{\ce{O-F}}
\end{mcq}
\why{Largest $\Delta$EN: O (3.5) vs H (2.1). O--F is small $\Delta$EN ---
neighbors.}
\distractor{D}{picks the two most electronegative atoms; polarity needs the
DIFFERENCE}

\begin{mcq}
  Solid magnesium conducts electricity, but solid \ce{MgCl2} does not. The
  best explanation is that
  \choices
    \correct{Mg has delocalized valence electrons; \ce{MgCl2} has ions
      fixed in a lattice}
    \choice{\ce{MgCl2} contains no charged particles}
    \choice{Mg atoms are smaller than \ce{Mg^2+} ions}
    \choice{\ce{MgCl2} has covalent bonds that trap electrons}
\end{mcq}
\why{Metallic sea vs immobile lattice ions.}
\distractor{B}{it contains ions --- they just can't move}

\begin{mcq}
  A compound melts at 801$^\circ$C, dissolves in water, conducts only when
  molten or dissolved, and shatters when struck. Its bonding is best
  described as
  \choices
    \choice{metallic}
    \choice{nonpolar covalent}
    \correct{ionic}
    \choice{polar covalent molecular}
\end{mcq}
\why{The full ionic property fingerprint.}

% ---- 2.2 PE curves ---------------------------------------------------------
\begin{mcq}
  On a potential-energy versus internuclear-distance curve for two atoms,
  the equilibrium bond length is found at
  \choices
    \choice{the distance where the energy equals zero}
    \correct{the distance at the minimum of the curve}
    \choice{the distance where attraction disappears}
    \choice{the largest distance plotted}
\end{mcq}
\why{The well minimum balances attraction and repulsion.}
\distractor{A}{zero energy = separated atoms, not the bond}

\begin{mcq}
  \ce{H2} has bond energy 432 kJ/mol and length 74 pm; \ce{F2} has bond
  energy 155 kJ/mol and length 142 pm. Compared with \ce{H2}, the PE curve
  for \ce{F2} has a minimum that is
  \choices
    \choice{deeper and at shorter distance}
    \choice{deeper and at longer distance}
    \choice{shallower and at shorter distance}
    \correct{shallower and at longer distance}
\end{mcq}
\why{Smaller bond energy = shallower; longer bond = minimum farther right.}

% ---- 2.3 lattice -----------------------------------------------------------
\begin{mcq}
  Which compound has the largest lattice energy magnitude?
  \choices
    \choice{\ce{NaCl}}
    \choice{\ce{KCl}}
    \correct{\ce{MgO}}
    \choice{\ce{CaS}}
\end{mcq}
\why{$2\times2$ charges AND the smallest such pair --- beats \ce{CaS} on
radius.}
\distractor{D}{right charges, larger ions}

\begin{mcq}
  Ionic solids are brittle because an applied stress
  \choices
    \choice{breaks the covalent bonds within each molecule}
    \correct{slides layers so that like charges align and repel}
    \choice{releases electrons from the lattice}
    \choice{converts the solid to its molten form}
\end{mcq}
\why{Layer shift $\to$ cation--cation/anion--anion contact $\to$ cleavage.}
\distractor{A}{no molecules and no covalent framework in an ionic solid}

\begin{mcq}
  \ce{LiF} and \ce{KBr} have ions of the same charges. \ce{LiF}'s melting
  point is much higher because
  \choices
    \correct{its smaller ions sit closer, strengthening the attraction}
    \choice{fluorine is more electronegative than bromine}
    \choice{\ce{LiF} has stronger covalent character}
    \choice{K and Br have more electrons to share}
\end{mcq}
\why{Same $|q_1q_2|$, smaller $r$ $\to$ larger Coulombic attraction.}
\distractor{B}{electronegativity is about sharing, not lattice distance}

% ---- 2.4 metals/alloys -----------------------------------------------------
\begin{mcq}
  Adding a small percentage of carbon to iron makes steel harder than pure
  iron because the carbon atoms
  \choices
    \choice{replace iron atoms in the lattice}
    \choice{form ionic bonds with iron}
    \correct{occupy interstitial holes and hinder layers from sliding}
    \choice{increase the number of delocalized electrons}
\end{mcq}
\why{Interstitial alloy: pinned planes resist deformation.}
\distractor{A}{that's a substitutional alloy (needs similar radii)}

\begin{mcq}
  Which property is common to BOTH metallic solids and ionic solids?
  \choices
    \choice{electrical conductivity in the solid state}
    \choice{malleability}
    \correct{attraction arising from electrostatic forces}
    \choice{discrete molecules in the lattice}
\end{mcq}
\why{Cation--sea and cation--anion attractions are both Coulombic.}

% ---- 2.5 Lewis -------------------------------------------------------------
\begin{mcq}
  The total number of valence electrons in the sulfite ion, \ce{SO3^2-}, is
  \choices
    \choice{24}
    \correct{26}
    \choice{25}
    \choice{32}
\end{mcq}
\why{$6 + 3(6) + 2 = 26$.}
\distractor{A}{forgot the charge}
\distractor{D}{counted 4 oxygens (sulfate)}

\begin{mcq}
  In the best Lewis diagram for \ce{HCN} (C central), the number of lone
  pairs on the nitrogen atom is
  \choices
    \correct{1}
    \choice{2}
    \choice{3}
    \choice{0}
\end{mcq}
\why{H--C$\equiv$N: triple bond + one lone pair completes N's octet.}

\begin{mcq}
  Which molecule has a central atom with LESS than an octet in its accepted
  Lewis diagram?
  \choices
    \choice{\ce{NH3}}
    \choice{\ce{H2O}}
    \correct{\ce{BF3}}
    \choice{\ce{CO2}}
\end{mcq}
\why{B holds 6; the double-bonded alternative puts + formal charge on F.}

% ---- 2.6 FC / resonance ----------------------------------------------------
\begin{mcq}
  In a Lewis diagram of \ce{CO2} drawn with one single bond and one triple
  bond, the formal charge on the single-bonded oxygen is
  \choices
    \correct{$-1$}
    \choice{$0$}
    \choice{$+1$}
    \choice{$-2$}
\end{mcq}
\why{$6 - 6 - 1 = -1$; this nonzero FC is why the symmetric two-double-bond
form is preferred.}

\begin{mcq}
  Experimental evidence that \ce{NO3-} is a resonance hybrid is that
  \choices
    \choice{it has 24 valence electrons}
    \choice{nitrogen carries a $+1$ formal charge}
    \correct{all three N--O bonds have identical, intermediate lengths}
    \choice{the ion has a trigonal planar shape}
\end{mcq}
\why{One short double + two long singles would show two lengths; hybrids show
one.}
\distractor{D}{shape is consistent with, but doesn't prove, resonance}

\begin{mcq}
  The carbonate ion \ce{CO3^2-} has three resonance forms. The carbon--oxygen
  bond order is
  \choices
    \choice{1}
    \correct{$4/3$}
    \choice{$3/2$}
    \choice{2}
\end{mcq}
\why{4 bonding pairs spread over 3 positions.}
\distractor{C}{that's \ce{NO2-}/\ce{O3} (2 positions)}

% ---- 2.7 VSEPR / hybrid / polarity ----------------------------------------
\begin{mcq}
  The molecular geometry of \ce{SF2} is
  \choices
    \choice{linear}
    \choice{trigonal planar}
    \correct{bent}
    \choice{tetrahedral}
\end{mcq}
\why{$6+14 = 20$ e$^-$: S has 2 bonds + 2 lone pairs = 4 domains, 2 LP $\to$
bent.}
\distractor{A}{ignores the lone pairs on S}

\begin{mcq}
  The F--N--F angle in \ce{NF3} is closest to
  \choices
    \choice{90$^\circ$}
    \correct{107$^\circ$}
    \choice{109.5$^\circ$}
    \choice{120$^\circ$}
\end{mcq}
\why{4 domains with 1 lone pair compressing the tetrahedral angle.}
\distractor{C}{forgets lone-pair compression}

\begin{mcq}
  The hybridization of the carbon atom in \ce{H2CO} (formaldehyde, C
  central) is
  \choices
    \choice{sp}
    \correct{sp$^2$}
    \choice{sp$^3$}
    \choice{sp$^3$d}
\end{mcq}
\why{3 domains (2 C--H, 1 C=O).}
\distractor{D}{not assessed on the AP exam and wrong anyway}

\begin{mcq}
  Which molecule is nonpolar despite containing polar bonds?
  \choices
    \choice{\ce{NH3}}
    \choice{\ce{H2O}}
    \choice{\ce{CHCl3}}
    \correct{\ce{CCl4}}
\end{mcq}
\why{Four equal dipoles in tetrahedral symmetry cancel.}
\distractor{C}{one H breaks the symmetry --- polar}

\examsection{II}{Free Response}{1 question \textbullet\ 20 minutes}{\calculatorok}

% ---- Long FRQ: CO2 vs SO2 + lattice ---------------------------------------
\frq{long}{10}{CED 2.1--2.7 synthesis \textbullet\ SP 1, 4, 6}

Carbon dioxide (\ce{CO2}) and sulfur dioxide (\ce{SO2}) are both triatomic
oxides, yet they differ in shape and polarity. Answer all parts.

\begin{parts}
  \item Draw the complete Lewis diagram of \ce{CO2}, and one complete
        resonance form of \ce{SO2}. Show all lone pairs.
        \workspace[3.4cm]
        \keyonly{\begin{apcanswer}\ce{CO2}: O=C=O, 2 lone pairs per O, none
        on C (16 e$^-$). \ce{SO2}: bent skeleton, one S=O and one S--O (or
        both drawn double --- accept), LONE PAIR ON S (18 e$^-$).\end{apcanswer}}
  \item Identify the molecular geometry and approximate bond angle of each
        molecule. \answerlines[2]{\ce{CO2}: linear, 180$^\circ$ (2 domains).
        \ce{SO2}: bent, slightly less than 120$^\circ$ (3 domains, 1 lone
        pair).}
  \item Explain why the two geometries differ, in terms of electron
        domains. \answerlines[2]{C has two bonding domains and no lone
        pairs; S has three domains including a lone pair, which occupies
        space and bends the O--S--O angle.}
  \item One of these molecules is polar. Identify it and justify using bond
        dipoles and geometry. \answerlines[2]{\ce{SO2} --- its bent shape
        prevents the two S--O dipoles from canceling, leaving a net dipole;
        \ce{CO2}'s linear geometry cancels its dipoles exactly.}
  \item The two S--O bonds in \ce{SO2} are found to be equal in length.
        Explain this observation. \answerlines[2]{\ce{SO2} is a resonance
        hybrid of the two equivalent forms; each bond has the same
        intermediate order, so equal lengths.}
  \item Solid \ce{CO2} (dry ice) sublimes at $-78^\circ$C, while \ce{MgO}
        melts near 2850$^\circ$C. Identify the particles occupying each
        solid's lattice and explain the enormous difference.
        \answerlines[3]{Dry ice: discrete \ce{CO2} molecules held by weak
        intermolecular attractions --- little energy to separate. \ce{MgO}:
        \ce{Mg^2+} and \ce{O^2-} ions in a lattice held by strong $2+/2-$
        Coulombic attraction throughout --- enormous energy to disrupt.}
\end{parts}

\begin{rubric}
  \rpoint{2}{(a) 1 pt \ce{CO2} correct (16 e$^-$, all lone pairs); 1 pt
    \ce{SO2} correct including the lone pair on S (18 e$^-$)}
  \rpoint{2}{(b) 1 pt per molecule: geometry AND angle}
  \rpoint{1}{(c) domain-count contrast, lone pair named as the bender}
  \rpoint{2}{(d) 1 pt identifies \ce{SO2}; 1 pt cancellation argument both
    directions}
  \rpoint{1}{(e) hybrid/average with equal bond order --- ``alternating''
    language earns 0}
  \rpoint{2}{(f) 1 pt particles named (molecules vs ions); 1 pt force-scale
    comparison tied to Coulomb}
\end{rubric}

\examtotal

\end{document}
